%-------------------------------------------------------------------------------
%	SECTION TITLE
%-------------------------------------------------------------------------------
\cvsection{Expérience}


%-------------------------------------------------------------------------------
%	CONTENT
%-------------------------------------------------------------------------------
\begin{cventries}
	\cventry
	{Ingénieur logiciel} % Intitulé du poste
	{Akkodis, en mission chez Thales} % Organisation
	{Gennevilliers, France} % Lieu
	{Octobre 2025 - Présent} % Dates
	{
		\begin{cvitems} % Description(s) des tâches/responsabilités
		\item{
			Développement d'un framework \textbf{C++} dédié à l'orchestration de réseaux (systèmes satellites).
		}
		\item{
			Collaboration au sein d'une équipe de 4 développeurs expérimentés.
		}
		\item{
			\textbf{C++} : Refonte complète d'une bibliothèque de sérialisation/désérialisation (\textbf{JSON}, \textbf{Protobuf}).
			\begin{itemize}
				\item{Utilisation avancée de la \textbf{métaprogrammation} (templates).}
				\item{Gestion de la représentation sérialisée des types primitifs (nombres flottants, signés, etc.).}
				\item{Génération de code \textbf{C++} à partir d'un DSL via \textbf{Xtend}.}
				\item{Tests unitaires \textbf{GoogleTest} et tests fonctionnels \textbf{Robot Framework}.}
			\end{itemize}
		}
		\item{
			Optimisation de l'environnement de développement et de la CI/CD :
			\begin{itemize}
				\item{Mise en place de \textbf{DevContainers} (VS Code)}
				\item{Maintenance et évolution des \textbf{Dockerfiles} et des pipelines \textbf{Jenkinsfiles}.}
			\end{itemize}
		}
		\item Développement backend (Python/FastAPI, Java/Spring) d'applications web de configuration réseau :
			\begin{itemize}
 				\item Code asynchrone (\textbf{async}, \textbf{await}), nouvelles fonctionnalités (endpoints REST) et refactorisation
				\item Écriture et maintenance des tests unitaires avec \textbf{pytest} et \textbf{JUnit}
				\item Évolution du simulateur d'équipements réseau pour faciliter les tests en environnement local
				\item Support technique aux équipes d'intégration sur les bancs d'essai et synchronisation avec les équipes de développement front-end
				\item Collaboration directe avec les ingénieurs système pour affiner les spécifications
			\end{itemize}
		\end{cvitems}
	}

	\cventry
	{Ingénieur logiciel embarqué} % Intitulé du poste
	{SII, en mission chez MBDA} % Organisation
	{Le Plessis-Robinson, France} % Lieu
	{Février 2024 - Octobre 2025} % Dates
	{
		\begin{cvitems} % Description(s) des tâches/responsabilités
		\item{
			Projet : Logiciels embarqués en \textbf{C++} et \textbf{Java (Swing)}, utilisant \textbf{IBM Rational Rhapsody} (Statecharts, \textbf{UML})
			\begin{itemize}
					\item{Réseau \textbf{Ethernet} et bus \textbf{1553}}
			\end{itemize}
		}
		\item{
			\textbf{Département Ingénierie Logicielle (\emph{SE})} : Membre d'une équipe de 4 développeurs, avec pour principales responsabilités :
			\begin{itemize}
				\item{\textbf{C++}, \textbf{Java} (\textbf{Swing} : interface graphique) : Conception, développement, refonte et correction de code \emph{legacy}}
				\item{Tests d'intégration des logiciels sur \textbf{bancs de test} : utilisation de moyens de test (simulateurs propriétaires) pour tester le système}
				\item{Génération personnalisée de \textbf{Linux embarqué} (type Yocto)}
				\item{Correction du \textbf{bug 2038} au niveau applicatif et au niveau du système d'exploitation Linux}
				\item{Drivers (modules noyau) Linux : recherche de nouveaux drivers ou patch de drivers \emph{legacy} pour un nouveau noyau Linux}
				\item{Migration vers une nouvelle version du compilateur (\textbf{gcc} 4 -> 12) : résolution des problèmes de compilation dus aux évolutions}
				\item{\textbf{Intégration} de code source \textbf{C++} fourni par une autre équipe : validation de l'impact du changement d'environnement}
			\end{itemize}
		}
		\item{
			Travail réalisé de manière autonome et à mon initiative, pour \textbf{améliorer les processus de travail et la qualité des livrables du projet} :
			\begin{itemize}
				\item{\textbf{Python}, \textbf{PowerShell}, \textbf{bash} et \textbf{Makefile} pour l'automatisation des tests, de la compilation et de la documentation}
				\item{Automatisation de tests d'une interface graphique \textbf{Java Swing} avec \emph{Robot} et \emph{java.lang.reflect}}
				\item{Automatisation de \textbf{IBM Rational Rhapsody} via l'API \textbf{Java}}
				\item{Automatisation de \textbf{Wireshark} via l'interface en ligne de commande \textbf{tshark}}
				\item{Écriture d'un générateur de code \textbf{C} en Python pour la sérialisation/désérialisation. \textbf{C++} : Génération de tests de données aléatoires pour le code généré}	
				\item{Configuration TCP/IP automatique \textbf{DHCP}/\textbf{bash}. Scripts \textbf{bash} très robustes pour le \textbf{NTP}}
			\end{itemize}
		}
		\end{cvitems}
	}

	\cventry
	{Ingénieur logiciel embarqué et simulation (contrat de thèse)} % Intitulé du poste
	{Laboratoire TIMA} % Organisation
	{Grenoble, France} % Lieu
	{Octobre 2022 - Janvier 2024} % Dates
	{
		\begin{cvitems} % Description(s) des tâches/responsabilités
		\item{
			Modélisation et simulation d'environnements de systèmes de contrôle :
			\begin{itemize}
				\item{Évaluation comparative de la simulation de systèmes cyber-physiques dans \textbf{Simulink} et \textbf{SystemC}}
			\end{itemize}
		}
		\item{
			Ingénierie et développement logiciel :
			\begin{itemize}
				\item{Modèles \textbf{Simulink/MATLAB} (blocs personnalisés utilisant \textbf{S-Function} en \textbf{C})}
				\item{Modèles \textbf{SystemC (C++)} (AMS/TLM) de SoC (\textbf{System-On-Chip})}
				\item{Restauration et refonte de code \emph{legacy} en \textbf{C++}, refonte de \textbf{Makefile}s}
			\end{itemize}
		}
		\end{cvitems}
	}

	\cventry
	{Enseignant} % Intitulé du poste
	{Université Grenoble Alpes, DLST} % Organisation
	{Grenoble, France} % Lieu
	{Janvier 2023 - Mai 2023} % Dates
	{
		\begin{cvitems} % Description(s) des tâches/responsabilités
		\item{Cours « Systèmes et environnement de programmation » pour étudiants de première année : travaux pratiques et exercices en classe}
		\item{\textbf{Scripts Bash}, bases de la programmation en \textbf{C} et modélisation simple par machines à états finis}
		\end{cvitems}
	}

	%---------------------------------------------------------
	\cventry
	{Stagiaire Master 2} % Intitulé du poste
	{Laboratoire TIMA} % Organisation
	{Grenoble, France} % Lieu
	{Février 2022 - Juin 2022} % Dates
	{
		\begin{cvitems} % Description(s) des tâches/responsabilités
		\item{Simulation parallèle de modèles de systèmes cyber-physiques/embarqués en \textbf{SystemC (C++)}}
		\item{Restauration et refonte de code \emph{legacy} \textbf{SystemC (C++)} et de \textbf{Makefile}s	}		
		\end{cvitems}
	}

	%---------------------------------------------------------
	\cventry
	{Stages} % Intitulé du poste
	{Kayentis, Givaudan, Abissa} % Organisation
	{Meylan, France, Vernier, Genève} % Lieu
	{Été 2021, 2019 et 2018} % Dates
	{
		\begin{cvitems} % Description(s) des tâches/responsabilités
		\item{Développement d'un script \textbf{PowerShell} pour un système de réponse automatique aux fautes}
		\item{Automatisation de navigateur web avec \textbf{Python Selenium} pour vérifier le fonctionnement d'une application Web}
		\item{Utilisation de \textbf{Jenkins} pour la surveillance automatique de l'application Web}
		\item{Recherche et déploiement d'une solution de surveillance de réseau \textbf{Ethernet} (Advanced Host Monitor)}
		\item{Recherche et déploiement d'une solution de gestion d'adresses IP (IPAM) (PhpIpam)}
		\end{cvitems}
	}

	%---------------------------------------------------------
\end{cventries}
